% CDI-PRBF Interface A: Visual Feedback Design Rationale
% HCI Academic Paper Style — Algorithmic & Feedback Logic Only
% No implementation details, no unrelated content.

\documentclass[11pt,a4paper]{article}

\usepackage[margin=1.2in]{geometry}
\usepackage{amsmath}
\usepackage{booktabs}
\usepackage{graphicx}
\usepackage{xcolor}
\usepackage{hyperref}
\usepackage{caption}
\usepackage{subcaption}
\usepackage{float}

\title{Interface A (Graphic Visualisation):\\
Algorithmic Rationale for Breathing-Guided Biofeedback}

\author{}
\date{}

\begin{document}
\maketitle

\section{Introduction}

This document describes the three-channel visual feedback logic of Interface~A in the
CDI-PRBF (Personalised Resonance Biofeedback) system. Interface~A adopts an
\textit{abstract-instrumental} representational style: a pulsing circle whose size,
colour, and ambient particle field each convey one independent dimension of the
user's physiological state. The same sensing pipeline also drives Interface~B
(metaphorical/ambient), making the two a fair pair for studying the effect of
visual framing on biofeedback efficacy.

The three feedback channels are:

\begin{enumerate}
  \item \textbf{Circle size} — breathing-guidance pacing (an exogenously
    generated signal, independent of physiology).
  \item \textbf{Circle colour} — real-time RSA (respiratory sinus arrhythmia)
    coherence between heart rate and the prescribed breathing pattern.
  \item \textbf{Particle burst} — a discrete reward event triggered when LF
    (low-frequency) heart-rate-variability power exceeds a threshold,
    signalling strong resonance.
\end{enumerate}

Each channel is described below in terms of its \textit{input signal},
\textit{transform}, and \textit{visual mapping}.

%-------------------------------------------------------------------------------
\section{Channel 1: Circle Size — Breathing Pacing}

\subsection{Input}
The pacing engine prescribes a target breathing frequency $f_{\text{bpm}}$
(breaths per minute) and an inhale-to-exhale ratio $R_I:R_E$. During the
frequency-sweep phase (Find PRBF), $f_{\text{bpm}}$ steps through ten discrete
values from 8.0 to 10.25~bpm; during the Control-Group phase it is fixed at
6.0~bpm (0.1~Hz).

\subsection{Transform}
Let the total cycle duration be
\[
T_{\text{cycle}} = \frac{60}{f_{\text{bpm}}} \quad\text{(seconds)}.
\]

The inhale and exhale segment durations are
\[
T_I = T_{\text{cycle}} \cdot \frac{R_I}{R_I + R_E}, \qquad
T_E = T_{\text{cycle}} \cdot \frac{R_E}{R_I + R_E}.
\]

Within each segment, phase progress $p \in [0,1]$ is defined as
\[
p(t) = \frac{t - t_{\text{seg-onset}}}{T_{\text{seg}}},
\]
clamped to $[0,1]$. The raw progress is reshaped by a quadratic ease-in-out
function to avoid linear-motion artefacts:
\[
\text{ease}(p) =
\begin{cases}
2p^2,                & p < 0.5, \\
-1 + (4 - 2p) \cdot p, & p \geq 0.5.
\end{cases}
\]

\subsection{Visual Mapping}
The circle's CSS \texttt{transform: scale} is linearly interpolated between
two anchor scales:
\[
S_{\text{inhale}} = 1.42, \qquad S_{\text{exhale}} = 0.72.
\]

During inhale, the circle expands:
\[
s(t) = (1 - \text{ease}(p)) \cdot S_{\text{exhale}} \;+\; \text{ease}(p) \cdot S_{\text{inhale}}.
\]

During exhale, the circle contracts:
\[
s(t) = (1 - \text{ease}(p)) \cdot S_{\text{inhale}} \;+\; \text{ease}(p) \cdot S_{\text{exhale}}.
\]

A concentric arc simultaneously drains clockwise, completing one full
circumference per cycle, providing a secondary spatial cue for cycle progress.

\paragraph{Rationale.}
The dynamic range $[0.72, 1.42]$ was chosen such that the circle is never
smaller than~72\,\% nor larger than~142\,\% of its rest size, keeping the
visual anchor recognisable across the full cycle. Quadratic easing mimics
natural diaphragmatic breathing kinematics (slow start/end, faster mid-phase),
which pilot testing suggested yields smoother perceived pacing than linear
interpolation.

%-------------------------------------------------------------------------------
\section{Channel 2: Circle Colour — RSA Coherence}

The colour channel operates in two regimes depending on the test mode.

\subsection{Find-PRBF Mode: Phase-Colour Mapping}
During the frequency sweep, colour serves as a simple phase indicator:
\begin{itemize}
  \item Inhale → Blue (\texttt{\#58a6ff}).
  \item Exhale → Green (\texttt{\#3fb950}).
\end{itemize}
No physiological feedback is encoded; the colour merely reinforces the
inhale/exhale transition, which is already signalled by size change.

\subsection{Control-Group Mode: RSA Coherence}

When the breathing frequency is held at 0.1~Hz, colour becomes a
physiologically driven channel. The system continuously computes the
cross-correlation between the user's instantaneous heart rate and the
prescribed breathing sine wave, yielding a real-time \textit{RSA coherence}
score.

\subsubsection{Signal Processing}

Let $\mathbf{h} = \{h_1, \dots, h_N\}$ be a sliding window of $N \leq 100$
heart-rate samples (BPM), each tagged with the breathing phase
$\phi_i \in [0,1]$ at the moment of measurement. The breathing reference
signal is synthesised from the animation phase:
\[
b_i = \sin(2\pi \cdot \phi_i), \quad i = 1 \dots N.
\]

Both signals are centred (zero-mean):
\[
\tilde{h}_i = h_i - \bar{h}, \qquad \tilde{b}_i = b_i - \bar{b}.
\]

The cross-correlation at lag $k$ is
\[
r(k) = \frac{\sum_{i} \tilde{h}_i \cdot \tilde{b}_{i-k}}
            {\sqrt{\sum_i \tilde{h}_i^2 \cdot \sum_i \tilde{b}_i^2}}.
\]

A lag search over $k \in [-15, +15]$ samples finds the peak absolute
correlation:
\[
C_{\text{raw}} = \max_{|k| \leq 15} \bigl| r(k) \bigr|,
\]
which is clamped to $[0, 1]$. The lag search accounts for individual
differences in the phase relationship between HR change and respiratory
phase (some individuals show HR peak slightly before or after end-inhale).

\subsubsection{Temporal Smoothing}

Raw coherence is smoothed by exponential moving average (EMA) with blend
factor $\alpha = 0.3$:
\[
C_{\text{smooth}}^{(t)} = \alpha \cdot C_{\text{raw}}^{(t)}
                         + (1 - \alpha) \cdot C_{\text{smooth}}^{(t-1)}.
\]

\subsubsection{Colour-Level Mapping with Hysteresis}

The smoothed coherence is discretised into four colour levels with
hysteresis to suppress flickering:

\begin{center}
\begin{tabular}{lll}
\toprule
\textbf{Level} & \textbf{Coherence Range} & \textbf{Colour} \\
\midrule
0 & $C < 0.4$                  & White \quad \texttt{rgb(255,255,255)} \\
1 & $0.4 \leq C < 0.6$        & Light Blue \quad \texttt{rgb(179,217,255)} \\
2 & $0.6 \leq C < 0.8$        & Sky Blue \quad \texttt{rgb(102,179,255)} \\
3 & $C \geq 0.8$              & Deep Blue \quad \texttt{rgb(26,128,230)} \\
\bottomrule
\end{tabular}
\end{center}

Transitions between levels require \textit{streak confirmation}: three
consecutive coherence values must lie in the candidate level's range
($\pm 0.1$ hysteresis band) before the colour changes. A 60-second
calibration delay after session start suppresses discolouration during
the initial settling period.

\subsubsection{Visual Rendering}

The mapped colour is applied simultaneously to: the circle's radial
gradient (lighter at top-left, fading to dark navy at the rim), an outer
ring stroke, the concentric arc, and the phase-text label. The circle
also carries a CSS \texttt{box-shadow} glow in the same colour. This
\textit{multi-element colour binding} ensures the chromatic signal is
perceptible even in peripheral vision.

\paragraph{Rationale.}
The four-level discretisation with hysteresis was chosen over a continuous
colour ramp because pilot participants reported that continuously morphing hues
were difficult to interpret as \textit{states}; the stepped palette with streak
confirmation provides clear, stable feedback about whether resonance quality is
improving, stable, or degrading. The calibration delay prevents early erratic
readings from anchoring an incorrect baseline impression.

%-------------------------------------------------------------------------------
\section{Channel 3: Particle Burst — LF-Power Reward}

\subsection{Input}
LF power (ms\textsuperscript{2}) is computed by the backend HRV analyser from
a sliding window of inter-beat intervals and streamed to the frontend at
approximately 1~Hz. LF power in the 0.04--0.15~Hz band is the primary HRV
marker of baroreflex-mediated resonance~\cite{lehrer2014}.

\subsection{Trigger Logic}

A discrete particle burst (``blossom'') fires when \textit{all} of the
following conditions are satisfied:

\begin{enumerate}
  \item \textbf{Threshold:} $\text{LF} \geq \tau$, where default
    $\tau = 500$~ms\textsuperscript{2} (user-configurable).
  \item \textbf{Streak:} Three consecutive readings above threshold
    (confirms the elevation is sustained, not a single-sample artefact).
  \item \textbf{Cooldown:} At least $T_{\text{cd}} = 120$~seconds since
    the previous burst.
  \item \textbf{Cap:} Fewer than $M = 5$ bursts have occurred in the
    current session.
\end{enumerate}

When all conditions hold, a burst is emitted and the streak counter resets.

\subsection{Progress Accumulator}

Between bursts, a continuous progress value $P \in [0,1]$ is computed as
\[
P = \min\!\left(1,\; \frac{t - t_{\text{last}}}{T_{\text{cd}}}\right)
      \cdot \left(0.5 + 0.5 \cdot \min\!\left(1,\; \frac{\text{LF}}{2\tau}\right)\right),
\]
where $t$ is the current session time and $t_{\text{last}}$ is the timestamp
of the most recent burst. The first factor models cooldown recovery; the second
factor modulates the rate by current LF power, so that higher LF power fills
the progress indicator faster (up to $2\times$ the baseline rate).

\subsection{Particle Animation}

Each burst spawns $N = 100$ particles from the screen centre. Each particle
$j$ is initialised as:

\begin{itemize}
  \item \textbf{Angle:} $\theta_j = \frac{2\pi j}{N} + \delta,\;
    \delta \sim \mathcal{U}(-0.3, 0.3)$ rad — near-uniform radial distribution
    with slight random jitter.
  \item \textbf{Speed:} $v_j \sim \mathcal{U}(120, 440)$ px/s.
  \item \textbf{Hue:} $H_j \sim \mathcal{U}(190, 270)$ — spanning cyan to
    purple in HSL colour space.
  \item \textbf{Life:} $L_j = 1.0$, decaying at
    $d_j \sim \mathcal{U}(0.4, 0.9)$ per second.
  \item \textbf{Size:} $r_j \sim \mathcal{U}(4, 12)$ px.
\end{itemize}

Particle dynamics follow a simple kinematic model with gravity:
\begin{align*}
  x_j(t+\Delta t) &= x_j(t) + v_{x,j} \cdot \Delta t, \\
  y_j(t+\Delta t) &= y_j(t) + v_{y,j} \cdot \Delta t, \\
  v_{y,j}(t+\Delta t) &= v_{y,j}(t) + 60 \cdot \Delta t \quad\text{(px/s\textsuperscript{2} downward)}.
\end{align*}

Rendering uses HSL colour with lightness proportional to remaining life
($55 + 20 \cdot L_j$), alpha faded by a factor of $1.2 \cdot L_j$, and a
12~px glow shadow. Particles are removed when $L_j \leq 0$; the animation
loop terminates when all particles are dead, clearing the overlay canvas.

\paragraph{Rationale.}
The burst is designed as a \textit{discrete positive reinforcer} for
achieving high LF power — a physiological state the user cannot directly
control but can influence through breathing quality. The three-condition
gate (threshold + streak + cooldown) prevents false triggering from
motion artefacts (which typically produce single-sample outliers) and
avoids satiation by limiting frequency (max $\approx 5$ per 10-minute
session). The particle aesthetic (cyan-to-purple with glow) was selected
to be visually distinct from the breathing circle's blue-green palette,
ensuring the reward event is perceptually salient without being
confusable with the continuous guidance channels.

%-------------------------------------------------------------------------------
\section{Summary of Signal-to-Channel Mapping}

\begin{table}[H]
\centering
\caption{Complete signal-to-channel mapping for Interface A (Graphic).}
\begin{tabular}{llll}
\toprule
\textbf{Signal} & \textbf{Source} & \textbf{Interface A} & \textbf{Interface B (for reference)} \\
\midrule
Breathing guidance  & Pacing engine          & Circle size (scale 0.72--1.42) & Cloud size \\
Resonance score     & RSA--guidance coherence& Circle colour (4-level blue scale) & Cloud colour \\
HRV reward          & LF power               & Particle burst (threshold-triggered) & Number of flowers \\
\bottomrule
\end{tabular}
\end{table}

All three visual channels operate simultaneously, are rendered on distinct
visual layers (CSS transform, CSS colour properties, HTML Canvas overlay), and
can be independently parameterised. The design isolates representational style
as the sole independent variable when compared with Interface~B — both
interfaces share identical signal processing, triggering logic, and update
rates.

%-------------------------------------------------------------------------------
\begin{thebibliography}{1}

\bibitem{lehrer2014}
P.M.\ Lehrer and R.\ Gevirtz,
``Heart rate variability biofeedback: how and why does it work?''
\textit{Frontiers in Psychology}, vol.~5, p.~756, 2014.

\bibitem{vaschillo2006}
E.\ Vaschillo, B.\ Vaschillo, and P.\ Lehrer,
``Characteristics of resonance in heart rate variability stimulated by biofeedback,''
\textit{Applied Psychophysiology and Biofeedback}, vol.~31, no.~2, pp.~129--142, 2006.

\bibitem{shaffer2017}
F.\ Shaffer and J.P.\ Ginsberg,
``An overview of heart rate variability metrics and norms,''
\textit{Frontiers in Public Health}, vol.~5, p.~258, 2017.

\end{thebibliography}

\end{document}
